\begin{proofbox}
\textbf{统一记号与目标。}
令 $P=(x,y)$，$y\ne0$，$c^2=a^2-b^2>0$，
$A_1=(-a,0),A_2=(a,0)$。利用 $x$ 轴对称性，
不妨先证 $y>0$；下半椭圆只需把 $y$ 换成 $|y|$。
证明的关键不是直接对角度求导，而是先证明
\[
\boxed{\tan(\alpha-90^\circ)=\frac{c^2|y|}{2ab^2}},\tag{$*$}
\]
再用 $0<|y|\le b$ 得到最大值。

\medskip
\textbf{换角公式补证：由 $\cot\alpha$ 到 $\tan(\alpha-90^\circ)$（两种推导）。}

这一步是前面证明用到的\textbf{三角恒等式补证}，不是替代向量法或几何法
对椭圆本身的证明。以下始终有 $90^\circ<\alpha<180^\circ$，
所以 $\sin\alpha>0$、$\tan\alpha<0$，函数均有定义。

\textbf{推导 A：正弦、余弦差角公式（直接且严格）。}
\begin{align*}
\tan(\alpha-90^\circ)
&=\frac{\sin(\alpha-90^\circ)}{\cos(\alpha-90^\circ)}\\
&=\frac{\sin\alpha\cos90^\circ-\cos\alpha\sin90^\circ}
        {\cos\alpha\cos90^\circ+\sin\alpha\sin90^\circ}\\
&=\frac{-\cos\alpha}{\sin\alpha}
=\boxed{-\cot\alpha}.
\end{align*}

\textbf{推导 B：正切差角公式 $+$ 极限（用主导项）。}
取 $\gamma\to90^{\circ-}$，令 $t=\tan\gamma\to+\infty$。
因 $\alpha-90^\circ\in(0,90^\circ)$，正切函数在该点连续，
所以
\begin{align*}
\tan(\alpha-90^\circ)
&=\lim_{\gamma\to90^{\circ-}}\tan(\alpha-\gamma)\\
&=\lim_{t\to+\infty}\frac{\tan\alpha-t}{1+t\tan\alpha}\\
&=\lim_{t\to+\infty}
\frac{\tan\alpha/t-1}{1/t+\tan\alpha}\\
&=-\frac{1}{\tan\alpha}=\boxed{-\cot\alpha}.
\end{align*}
本题 $\tan\alpha<0$，故分子和分母的主导项
分别为 $-t$ 和 $t\tan\alpha$，二者同趋负无穷，
但比值趋于 $-1/\tan\alpha>0$。注意：
$\tan90^\circ$ \textbf{没有定义}；“无穷大比较”仅是上述\textbf{极限运算}的直观简称。

\textbf{接回椭圆：}由向量法或余弦定理得到
$\displaystyle\cot\alpha=-\frac{c^2|y|}{2ab^2}$ 后，
便可严格推出
\[
\tan(\alpha-90^\circ)=-\cot\alpha
=\frac{c^2|y|}{2ab^2}.
\]
\textbf{高考书写建议：}正式解答优先用推导 A（也可直接引用恒等式）；
推导 B 主要帮助理解为何不能把未定义的 $\tan90^\circ$ 生硬代入。

\medskip
\textbf{证法一：向量点积 $+$ 行列式（最短、最稳的代数证明）。}

从 $P$ 出发的两个向量为
\[
\vec u=\overrightarrow{PA_1}=(-a-x,-y),\quad
\vec v=\overrightarrow{PA_2}=(a-x,-y).
\]
它们的点积与平面行列式分别是
\[
\vec u\cdot\vec v=x^2+y^2-a^2
=-\frac{c^2}{b^2}y^2<0,
\]
\[
|\det(\vec u,\vec v)|=2a|y|>0.
\]
所以夹角是钝角，而且
\[
\cot\alpha=\frac{\vec u\cdot\vec v}
{|\det(\vec u,\vec v)|}
=-\frac{c^2|y|}{2ab^2}.
\]
由于 $\alpha=90^\circ+\delta$、$\delta\in(0,90^\circ)$，
$\cot\alpha=-\tan\delta$，于是得到 $(*)$。
显然 $|y|\le b$，等号仅在 $x=0$ 时成立，故
\[
\alpha\le90^\circ+\arctan\frac{c^2}{2ab},
\]
且等号仅在 $P=(0,\pm b)$ 时成立。

\medskip
\textbf{证法二：余弦定理 $+$ 面积公式（不需要向量）。}

设 $r_1=PA_1$、$r_2=PA_2$，底边 $A_1A_2=2a$。
由余弦定理
\[
\cos\alpha=\frac{r_1^2+r_2^2-4a^2}{2r_1r_2}.
\]
注意
\[
r_1^2+r_2^2=2(x^2+y^2+a^2),
\]
故
\[
\cos\alpha=\frac{x^2+y^2-a^2}{r_1r_2}
=-\frac{c^2y^2}{b^2r_1r_2}<0.
\]
另一方面，用底边与高求面积，再用两边夹角求面积：
\[
a|y|=S_{\triangle A_1PA_2}
=\frac12r_1r_2\sin\alpha,
\qquad \sin\alpha=\frac{2a|y|}{r_1r_2}.
\]
两式相除，得到
\[
\cot\alpha=\frac{\cos\alpha}{\sin\alpha}
=-\frac{c^2|y|}{2ab^2}.
\]
随后利用 $|y|\le b$ 即证最大值。此法强调：
\textbf{面积给正弦，余弦定理给余弦；一除便得到角度。}

\medskip
\textbf{证法三：底角 $+$ 正切加法公式（初等几何与三角恒等式）。}

仍设 $y>0$。三角形 $A_1PA_2$ 的两个底角分别为
$\beta_1=\angle PA_1A_2$、$\beta_2=\angle PA_2A_1$，则
\[
\tan\beta_1=\frac{y}{x+a},\qquad
\tan\beta_2=\frac{y}{a-x}.
\]
由椭圆方程
\[
\tan\beta_1\tan\beta_2
=\frac{y^2}{a^2-x^2}=\frac{b^2}{a^2}<1,
\]
因此 $0<\beta_1+\beta_2<90^\circ$，这也说明
$\alpha=180^\circ-(\beta_1+\beta_2)>90^\circ$。
正切加法公式给出
\begin{align*}
\tan(\beta_1+\beta_2)
&=\frac{\frac{y}{x+a}+\frac{y}{a-x}}
{1-\frac{y^2}{a^2-x^2}}\\
&=\frac{2ay}{a^2-x^2}\cdot\frac{a^2}{a^2-b^2}
=\frac{2ab^2}{c^2y}.
\end{align*}
故
\[
\tan\alpha=-\frac{2ab^2}{c^2y},\qquad
\tan(\alpha-90^\circ)=\frac{c^2y}{2ab^2}.
\]
式中\textbf{不能}把 $\alpha$ 与 $\beta_1+\beta_2$ 当成同一个角。

\medskip
\textbf{证法四：圆周角与外接圆圆心（最有几何味道）。}

对上半椭圆的 $P$ 作 $\triangle A_1PA_2$ 外接圆。
因 $A_1,A_2$ 关于 $y$ 轴对称，圆心必是 $Q=(0,h)$。
利用 $QA_2=QP$，得
\[
a^2+h^2=x^2+(y-h)^2
\quad\Longrightarrow\quad
h=\frac{x^2+y^2-a^2}{2y}
=-\frac{c^2y}{2b^2}<0.
\]
设 $\varphi=\arctan(|h|/a)$，由等腰三角形
$QA_1A_2$，其\textbf{较小圆心角}为
$\angle A_1QA_2=180^\circ-2\varphi$。
$P$ 在 $A_1A_2$ 的上方小弧上，$\alpha$ 所对的是下方\textbf{大弧}，
由圆周角定理
\[
\alpha=\frac12\bigl(360^\circ-\angle A_1QA_2\bigr)
=90^\circ+\varphi
=90^\circ+\arctan\frac{c^2y}{2ab^2}.
\]
圆心 $Q$ 越低，$|h|$ 越大，顶角越大；
$y\le b$，故最大值只能在短轴端点取得。

\medskip
\textbf{证法五：三角参数方程（把最值变成 $|\sin t|\le1$）。}

设 $P=(a\cos t,b\sin t)$，排除 $\sin t=0$（长轴端点）。
由
\[
x^2+y^2-a^2
=a^2\cos^2t+b^2\sin^2t-a^2
=-c^2\sin^2t,
\]
而两向量的行列式绝对值为 $2ab|\sin t|$，有
\[
\cot\alpha=-\frac{c^2}{2ab}|\sin t|,
\qquad
\alpha=90^\circ+\arctan\biggl(\frac{c^2}{2ab}|\sin t|\biggr).
\]
由 $0<|\sin t|\le1$，当且仅当
$|\sin t|=1$，即 $t=\pi/2$ 或 $3\pi/2$ 时最大，
对应 $(0,\pm b)$。没有求导，也没有对斜率分类讨论。

\medskip
\textbf{证法六：原稿的“斜率乘积定值”法（校正符号和角的定义）。}

设
\[
k_1=\frac{y}{x+a},\qquad k_2=\frac{y}{x-a},
\qquad k_1k_2=\frac{y^2}{x^2-a^2}=-\frac{b^2}{a^2}.
\]
当 $y>0$ 时，$k_1>0,k_2<0$。设两个\textbf{底角}
$\beta_1=\arctan k_1$、$\beta_2=\arctan(-k_2)$，则
$\alpha=180^\circ-(\beta_1+\beta_2)$。
于是
\[
\tan\alpha=-\frac{k_1-k_2}{1+k_1k_2}
=-\frac{2ab^2}{c^2y}.
\]
下半椭圆用 $|y|$ 对称延拓。特别提醒：
$\tan\alpha_{\max}=-2ab/c^2$，
绝不能只对其正切取绝对值后，把锐角误作最大顶角。

\medskip
\textbf{最后一步：短轴端点直接计算最大角（纯直角三角形收尾）。}

由前六种证法，最大值仅在 $P=(0,\pm b)$ 处发生。
以 $P=(0,b)$ 为例，设 $O=(0,0)$，
$PO$ 平分顶角且 $PO\perp A_1A_2$，所以
\[
\tan\frac{\alpha_{\max}}2=\frac{OA_2}{PO}=\frac ab
\quad\Longrightarrow\quad
\boxed{\alpha_{\max}=2\arctan\frac ab}.
\]
由正切倍角公式（注意 $\alpha_{\max}>90^\circ$）
\[
\tan\alpha_{\max}
=\frac{2(a/b)}{1-(a/b)^2}
=\boxed{-\frac{2ab}{a^2-b^2}}.
\]
由于 $y\ne0$，恒有 $\alpha>90^\circ$；
当 $P\to(\pm a,0)$ 时 $|y|\to0$，$\alpha\to90^\circ$，
但端点处角度无定义，所以\textbf{只有下确界而无最小值}。

\medskip
\textbf{延伸证明：等角 $\Longleftrightarrow$ 等面积 $\Longleftrightarrow$ 等高。}

由 $(*)$，固定椭圆参数 $a,b$ 时，$\alpha$ 与 $|y|$ 严格一一对应。
再利用
\[
S=\frac12\cdot 2a\cdot|y|=a|y|,
\qquad |x|=a\sqrt{1-\frac{y^2}{b^2}},
\]
得三者互相决定，且角度只能确定 $|x|,|y|$，
一般有四个轴对称位置。

\medskip
\textbf{关联 C003：两种角度为什么会同时在短轴最大？}

设 $\theta=\angle F_1PF_2$，焦点 $F_1(-c,0),F_2(c,0)$。
焦点三角形有 $PF_1+PF_2=2a$；由余弦定理
\[
4c^2=4a^2-4(PF_1)(PF_2)\cos^2\frac\theta2
\quad\Rightarrow\quad
(PF_1)(PF_2)=\frac{b^2}{\cos^2(\theta/2)}.
\]
结合面积
$\frac12(PF_1)(PF_2)\sin\theta=c|y|$，得
\[
\tan\frac\theta2=\frac{c|y|}{b^2}.
\]
与顶角公式 $(*)$ 联立：
\[
\boxed{\tan(\alpha-90^\circ)=\frac{c}{2a}\tan\frac\theta2}.
\]
故二者都随 $|y|$ 增大而增大，但取值范围与三角形含义不同。
\end{proofbox}
